% =====================================================================================================
\section{External databases and how the data are fetched}\label{sec:acquisition}
% =====================================================================================================

All inputs come from public services and are obtained by the scripts, with no manual downloads except the NICER CALDB
archive (\cref{sec:get-nicer}). \Cref{tab:sources} lists every external database, what is taken from it and how.
\Cref{fig:acq} shows how requests flow from the scripts to local storage and to the download log.

{\small\setlength{\tabcolsep}{3pt}
\begin{longtable}{@{}>{\raggedright\arraybackslash}p{3.2cm}>{\raggedright\arraybackslash}p{4.7cm}>{\raggedright\arraybackslash}p{4.4cm}>{\raggedright\arraybackslash}p{3.6cm}@{}}
\caption{External databases used by the workflow. ``Records'' and volumes are from \file{logs/downloads.jsonl}
(\cref{sec:provenance}).}\label{tab:sources}\\
\toprule
Database & What is taken & Access & Code \\
\midrule\endfirsthead
\toprule Database & What is taken & Access & Code \\ \midrule\endhead
\bottomrule\endfoot
\multicolumn{4}{@{}l}{\emph{Labels and positions}}\\
BlackCAT \citep{CorralSantana2016} & dynamically confirmed BHs (v1--v4); candidate list (v3c, descriptive) & names transcribed into \file{config.py}; candidates checked against the BlackCAT web pages & \file{config.py} (v2, v3c blocks) \\
\citet{Marcel2026}, Table 1 & 25 dynamically confirmed BHs (v7 onwards) & names transcribed into the scripts & \file{38_v7c2_select.py}, \file{47_v9_select.py} \\
\citet{Galloway2008} & type-I bursters, appendix A (v1--v3), Table 3 (v4b2) & names transcribed into \file{config.py} & \file{config.py}, \file{20_v4b2_select.py} \\
MINBAR DR1 \citep{Galloway2020} & burst list, observation list, source table with positions & figshare files 23201936, 24131849, 22851017 (2.5\,MB, 51.9\,MB, 0.7\,MB), MD5 checked & \file{31_v7b1_minbar.py}, \code{v7lib.load_minbar} \\
Pulsar lists & AMXPs and slow pulsars (Patruno \& Watts 2021); HMXB/LMXB catalogues of Liu et al. & names transcribed; catalogue tables from CDS (8 files) & \file{28_v6_select.py}, \file{38_v7c2_select.py} \\
CDS Sesame (SIMBAD) & RA/Dec of every named source & \code{GET https://cds.unistra.fr/cgi-bin/nph-sesame/-oI/SNV?<name>}; the \texttt{\%J} line is parsed; results cached in CSV & \code{sesame()} in \file{20}, \file{28}, \file{38} \\
\multicolumn{4}{@{}l}{\emph{Observation catalogues}}\\
RXTE MissionLongData & per-source FITS observation list: \code{OBSID}, \code{TIME}, \code{EXPOSURE}, \code{STD1RATE}, \code{RA}, \code{DEC}, \code{NPCU} & directory listing and \file{<name>.fits.gz} under \file{.../xte/data/archive/MissionLongData/} & \file{02}, \file{03}, \file{20}, \file{35}, \file{41} \\
HEASARC TAP \code{xtemaster} & RXTE pointings near a position (sources without a MissionLongData file) & ADQL cone query, HTTP POST to \file{.../xamin/vo/tap/sync} & \file{41_v8b_select.py} \\
HEASARC TAP \code{nicermastr} & NICER observations near a position: \code{obsid}, \code{time}, \code{exposure}, \code{num_fpm} & same service; one query per source, cached per source & \file{47_v9_select.py} \\
\multicolumn{4}{@{}l}{\emph{Data archives}}\\
HEASARC RXTE archive & Standard-2 Standard Products, Standard-1 light curves, event files, HEXTE products & HTTPS file tree; Python \code{urllib} (event files: \code{curl}); 46\,731 records, 5.2\,GB & \code{common.download}, \code{v8lib.download_curl} \\
HEASARC NICER (AWS open-data mirror) & cleaned (\code{_cl}) and unfiltered (\code{_ufa}) event files & \code{curl} HTTP range requests; 3\,384 records, 229.5\,GB read & \file{v9lib.py}, \file{v10lib.py}, \file{v11lib.py}, \file{v12_bkg3c50.sh} \\
MAXI/GSC & source index; 1-day light curves (2--4, 4--10, 10--20\,keV) & 13 category pages \file{http://maxi.riken.jp/top/lc_<cat>.html}; \file{http://maxi.riken.jp/star_data/<J>/<J>_g_lc_1day_all.dat} (143 files) & \file{38_v7c2_select.py} \\
GitHub (de Beurs et al.) & processed MAXI data and code of \citet{deBeurs2022} & GitHub API for the commit and file tree, then raw files at commit \texttt{488b785}; 487 files & \file{37_v7c1_repro.py} \\
\multicolumn{4}{@{}l}{\emph{Calibration}}\\
RXTE responses & \file{xp<obs>.rsp} per observation & with the Standard Products & \code{spectra.read_rsp} \\
NICER CALDB & 3C50 background libraries & \file{goodfiles_nicer_xti.tar.gz} (112\,180\,364 bytes, SHA256 recorded in the decision log) & \file{v12_bkg3c50.sh} \\
HI4PI & Galactic $N_\mathrm{H}$ at MAXI source positions & HEASoft \code{nh} with its HI4PI map & \file{v8d_heasoft.sh} \\
\end{longtable}}

\begin{figure}[htbp]
\centering
\begin{tikzpicture}[x=1cm, y=1cm,
  h/.style={fc data, text width=3.3cm, minimum height=12mm},
  m/.style={fc proc, text width=4.3cm, minimum height=12mm},
  s/.style={fc, draw=inkgrey, fill=boxgrey, text width=4.2cm, minimum height=12mm}]
  \node[fc head] at (1.75,0.9) {service};
  \node[fc head] at (7.0,0.9) {access method (code)};
  \node[fc head] at (12.75,0.9) {local result};
  \node[h] (h1) at (1.75,0)    {HEASARC RXTE file tree};
  \node[h] (h2) at (1.75,-1.6) {HEASARC TAP (xamin)};
  \node[h] (h3) at (1.75,-3.2) {NICER AWS mirror};
  \node[h] (h4) at (1.75,-4.8) {CDS Sesame; MINBAR (figshare); MAXI; GitHub};
  \node[m] (m1) at (7.0,0)    {\code{common.download}: skip if present, 3 tries, no retry on 404, write \file{.part} then rename};
  \node[m] (m2) at (7.0,-1.6) {ADQL cone query by POST, 4 tries; VOTable parsed with \code{astropy}};
  \node[m] (m3) at (7.0,-3.2) {\code{curl} HEAD, then range requests (16 or 64\,MB), 4 tries each; incremental gzip/FITS parser};
  \node[m] (m4) at (7.0,-4.8) {\code{common.download} or \code{urllib}; Sesame answers cached};
  \node[s] (s1) at (12.75,0)    {\file{data/raw/...} or \file{XRB_EXTERNAL_RAW/...}};
  \node[s] (s2) at (12.75,-1.6) {\file{data/raw/catalogs/nicer/nicermastr_<source>.csv}};
  \node[s] (s3) at (12.75,-3.2) {prefix and tail files (v9, v10) or memory only (v11, v13)};
  \node[s] (s4) at (12.75,-4.8) {\file{data/raw/minbar/}, \file{data/v7c/}, caches};
  \foreach \i in {1,2,3,4} {\draw[fc arr] (h\i) -- (m\i); \draw[fc arr] (m\i) -- (s\i);}
  \node[fc, draw=tagpost, dashed, fill=white, text width=15.2cm, minimum height=7mm, anchor=north west] at (0.1,-5.75)
    {every downloaded or streamed file: one JSON line in \file{logs/downloads.jsonl} with \code{url}, \code{path}, \code{bytes}, \code{sha256}, \code{utc} (and \code{tool})};
\end{tikzpicture}
\caption{Data acquisition. Each service is reached through one access function. Downloads are idempotent: a file that
already exists is not fetched again, and every new file is logged with its checksum.}\label{fig:acq}
\end{figure}

% -----------------------------------------------------------------------------------------------------
\subsection{Label lists and positions}\label{sec:get-labels}
Labels are attached to \emph{sources}, never to observations. The label sources are published tables, and the realised
source lists are the \file{data/*/sources.csv} files:
\begin{itemize}
\item \textbf{BHs, v1--v4:} the dynamically confirmed systems of BlackCAT Table~4, transcribed into the v2 block of
\file{config.py} with their MissionLongData catalogue names.
\item \textbf{BHs, v7 onwards:} the 25 systems of \citet{Marcel2026} Table~1, written as a list in the selection scripts.
\item \textbf{NSs:} type-I bursters of \citet{Galloway2008} (appendix A sections cited per source in \file{config.py}). From
v7 the MINBAR source table (115 bursters) is used instead; the external RXTE set adds accretion-powered pulsars.
\item \textbf{BH candidates} (BlackCAT, X-ray evidence only) are never labels. In v3c they were scored as prediction
targets, after the list had been checked against the BlackCAT web pages and frozen (\file{reports/v3c_candidate_verification.md}).
\end{itemize}

Positions come from three places, depending on the round:
\begin{itemize}
\item \textbf{v1--v3:} the median pointing in the source's RXTE catalogue.
\item \textbf{v4b2 onwards:} CDS Sesame. The script requests \code{https://cds.unistra.fr/cgi-bin/nph-sesame/-oI/SNV?<name>},
reads RA and Dec from the line starting with \texttt{\%J}, and caches the answer (\file{data/v7c/sesame_positions.csv},
\file{data/v4b2/missionlongdata_positions.csv}).
\item \textbf{NSs from v7:} the MINBAR source table (\file{minbar_sources_v2.7.fits}, columns \code{NAME}, \code{RA_OBJ},
\code{DEC_OBJ}).
\end{itemize}
The MINBAR burst table (\file{minbar.txt}, machine-readable format) supplies burst start times and durations
(\code{v7lib.load_minbar}). The three MINBAR files are downloaded from figshare; the MD5 sums of the two burst-related files
are fixed in \code{V7_MINBAR_FILES} in \file{config.py}.

% -----------------------------------------------------------------------------------------------------
\subsection{RXTE: catalogues and archive files}\label{sec:get-rxte}
\paragraph{Observation catalogues.}
The archive root is \file{https://heasarc.gsfc.nasa.gov/FTP/xte/data/archive/}. Its \file{MissionLongData/} directory
holds one FITS table per catalogued target. \file{02_catalogs.py} and \file{03_select_observations.py} download
\file{MissionLongData/<name>.fits.gz} and read extension 1 with \code{astropy}. The observation time is
\code{TIME + MJDREFI + MJDREFF} (MJD), and duplicated \code{OBSID}s are dropped. For v4b2, the whole directory listing
(291 catalogues) was downloaded and each catalogue's median pointing was matched to Sesame positions within $0.1^\circ$.
For new BHs without a catalogue file (v8b), the HEASARC TAP service is queried instead:
\begin{verbatim}
SELECT obsid, target_name, ra, dec, time, exposure FROM xtemaster
WHERE CONTAINS(POINT('ICRS',ra,dec),CIRCLE('ICRS',<ra>,<dec>,0.1))=1
\end{verbatim}
The query is sent by HTTP POST to \file{https://heasarc.gsfc.nasa.gov/xamin/vo/tap/sync} with \code{REQUEST=doQuery},
\code{LANG=ADQL}.

\paragraph{Archive layout and file names.} An ObsID \texttt{PPPPP-TT-VV-SS} lives under
\file{AO<n>/P<PPPPP>/<obsid>/}. The files used are listed below.

\begin{table}[htbp]
\centering\small
\caption{RXTE files read per observation (\texttt{o} = ObsID without hyphens).}\label{tab:rxtefiles}
\begin{tabularx}{\textwidth}{@{}>{\raggedright\arraybackslash}p{5.3cm} >{\raggedright\arraybackslash}X l@{}}
\toprule
Path below the ObsID directory & Content & Used in \\
\midrule
\file{stdprod/xp<o>_s2.pha.gz} & PCA Standard-2 total (source + background) spectrum, 129 channels & v1--v8 \\
\file{stdprod/xp<o>_b2.pha.gz} & \code{pcabackest} model background spectrum, same PCUs and GTIs & v1--v8 \\
\file{stdprod/xp<o>.rsp.gz} & response (RMF$\times$ARF) for that PCU combination, with \code{EBOUNDS} & v1--v8 \\
\file{pca/FS46_*.gz} & Standard-1: 0.125\,s counts per PCU (\code{XeCntPcu0}--\code{4}) and array rates & dead time (v2+), timing (v5+) \\
\file{pca/SE*.evt.gz} & event-mode files (\code{TIME}, \code{PCUID}) & v8c \\
\file{stdprod/xh<o>_s1.pha}, \file{_b1.pha} + RMF/ARF & HEXTE cluster-B spectra and responses & v4b3 \\
\bottomrule
\end{tabularx}
\end{table}

\paragraph{Locating the AO directory.} The AO number is not stored in the catalogue. \file{04_fetch_process.py} guesses
it from the proposal prefix (\code{5}$\to$AO5, \dots, \code{90}$\to$AO9, \code{91}$\to$AO10, \dots, \code{96}$\to$AO15). If the
guess returns HTTP 404, it tries AO4--AO16 in turn and caches the AO that worked for the proposal. (The v4b2 copy of the
script also scans AO1--AO4 for the early gain epochs.) The Standard-1 file names are taken from the HTML index of the
\file{pca/} directory (regular expression \verb|FS46_[^"]+\.gz|).

\paragraph{Download function.} \code{common.download(url, relative_path)} does the following:
\begin{enumerate}
\item returns the local file if it exists (re-runs never re-download);
\item otherwise fetches the URL with \code{urllib} (45\,s timeout), retrying up to three times with a 2\,s pause, but not
on HTTP 404;
\item writes the data to \file{<name>.part} and renames it, so an interrupted run leaves no truncated files;
\item appends a record with \code{url}, \code{path}, \code{bytes}, \code{sha256} and \code{utc} to
\file{logs/downloads.jsonl}, under a lock because downloads run in a thread pool (4 workers in
\file{04_fetch_process.py}).
\end{enumerate}
The RXTE event files of v8c (2.4\,GB) were fetched with \code{v8lib.download_curl}, which streams with \code{curl} into a
\file{.part} file and writes the same log record.

% -----------------------------------------------------------------------------------------------------
\subsection{NICER: TAP query, mirror and range requests}\label{sec:get-nicer}
\paragraph{Observation list.} For every source, \file{47_v9_select.py} queries the HEASARC TAP service:
\begin{verbatim}
SELECT obsid, name, ra, dec, time, exposure, num_fpm FROM nicermastr
WHERE CONTAINS(POINT('ICRS',ra,dec),CIRCLE('ICRS',<ra>,<dec>,0.05))=1
\end{verbatim}
The query runs with 6 parallel requests and 4 tries each. The VOTable is saved per source as
\file{data/raw/catalogs/nicer/nicermastr_<source>.csv}. The later rounds (v11, v13) select from these cached tables, so
all three NICER observation sets come from the same catalogue snapshot.

\paragraph{File location.} Cleaned event files are read from HEASARC's AWS open-data mirror, which is byte-identical to
the HEASARC site (checked on the first 4\,MB) and about ten times faster from the analysis machine:
\begin{center}
\file{https://nasa-heasarc.s3.amazonaws.com/nicer/data/obs/<YYYY_MM>/<obsid>/xti/event_cl/ni<obsid>_0mpu7_cl.evt.gz}
\end{center}
\code{YYYY_MM} is the month of the catalogue start time. If the file is not found (HTTP 403 or 404 on the mirror), the
previous and next months are tried. The unfiltered file \file{..._0mpu7_ufa.evt.gz} is fetched only for the 3C50
background (\cref{sec:3c50}).

\paragraph{Range requests and the streaming parser.} The files are gzip-compressed FITS tables of up to several GB, so
they are never downloaded blindly:
\begin{enumerate}
\item A HEAD request (\code{curl -I}) checks that the file exists and gives its size.
\item The file is fetched in HTTP range requests (\code{curl -r a-b}; options \code{--connect-timeout 30},
\code{--speed-limit 20000}, \code{--speed-time 60}), each with up to 4 tries and a 5\,s pause.
\item Each piece is passed to an incremental parser (\code{v9lib.PrefixParser}). The parser decompresses the gzip stream
with \code{zlib}, reads the primary and \code{EVENTS} headers from 2880-byte blocks, and converts complete binary-table
rows to arrays of \code{TIME}, \code{PI} and \code{DET_ID} using a numpy dtype built from the \code{TFORMn} keywords.
\end{enumerate}
No HEASoft is needed to read the events. The three reading modes are listed in \cref{tab:nicermodes}. Parsing was checked
to agree with \code{astropy} and across the modes (implementation checks of v9--v11).

\begin{table}[htbp]
\centering\small
\caption{NICER reading modes.}\label{tab:nicermodes}
\begin{tabularx}{\textwidth}{@{}>{\raggedright\arraybackslash}p{1.9cm} >{\raggedright\arraybackslash}p{3.1cm} >{\raggedright\arraybackslash}p{1.9cm} >{\raggedright\arraybackslash}X >{\raggedright\arraybackslash}X@{}}
\toprule
Round & Function & Request size & Stop rule & Storage \\
\midrule
v9 (prefix) & \code{v9lib.stream_prefix} & 16\,MB & $\ge6$ gap-free 128\,s segments, end of file, or 300\,MB & compressed prefix + JSON metadata in \file{XRB_EXTERNAL_RAW/nicer/<obsid>/} \\
v10 (tail) & \code{v10lib.download_tail} & 64\,MB & end of file or 1\,GB compressed & tail in \file{nicer_tail/<obsid>/}; prefix + tail parsed as one stream \\
v11, v13 (stream) & \code{v11lib.stream_events} & 64\,MB, 6 in parallel & end of file or 1\,GB & nothing stored (SHA256 of the bytes read is logged); v13 keeps cleaned files of faint observations for one check \\
\bottomrule
\end{tabularx}
\end{table}

\paragraph{3C50 inputs and the CALDB.} For the background model, the script \file{v12_bkg3c50.sh} runs inside WSL Ubuntu
with HEASoft 6.37.1. It:
\begin{enumerate}
\item downloads the cleaned and unfiltered event files with \code{curl} (3 tries) and prints their SHA256 and size, which
the Python driver writes to the download log;
\item unzips them and runs \code{nibackgen3C50} (\cref{sec:3c50});
\item keeps only the resulting total and background spectra (\file{tot_<obsid>.pi}, \file{bkg_<obsid>.pi});
\item deletes the event files.
\end{enumerate}
The NICER CALDB (\file{goodfiles_nicer_xti.tar.gz} with \file{caldb.config} and \file{alias_config.fits}) was downloaded
once before the v12 data and unpacked under \file{~/xrb-pilot-data/caldb}; its size and SHA256 are in the decision log.

% -----------------------------------------------------------------------------------------------------
\subsection{MAXI, published code and absorption}\label{sec:get-maxi}
\begin{itemize}
\item \textbf{MAXI matching.} The label lists are frozen first. The MAXI source index is read from the 13 category pages
(\file{http://maxi.riken.jp/top/lc_<cat>.html}). MAXI names are resolved with Sesame and matched to the label sources within
$0.1^\circ$, with a $0.3^\circ$ fallback to the position encoded in the MAXI J-name. MAXI sources whose name contains
``with'' (blends) are excluded, as are ambiguous matches (\cref{sec:maxi}).
\item \textbf{MAXI light curves.} The 1-day light curves are downloaded from
\file{http://maxi.riken.jp/star_data/<J>/<J>_g_lc_1day_all.dat} with \code{common.download}.
\item \textbf{Published code.} For the reproduction of \citet{deBeurs2022}, \file{37_v7c1_repro.py} asks the GitHub API
for the current commit and file tree of \file{zdebeurs/3ML_methods_for_XRB_classification}. It then downloads every file
from \file{raw.githubusercontent.com} at that commit (\texttt{488b785}).
\item \textbf{Absorption.} For the absorption correction, \file{v8d_heasoft.sh} runs HEASoft \code{nh} (HI4PI, average
within $0.1^\circ$) at each MAXI source position. It also runs XSPEC to tabulate \code{tbabs}$\times$power-law
transmissions (\cref{sec:maxi}).
\end{itemize}

% -----------------------------------------------------------------------------------------------------
\subsection{Provenance log and verification}\label{sec:provenance}
\file{logs/downloads.jsonl} has one JSON object per downloaded or streamed file. Its fields are \code{url}, \code{path}
(or a note such as ``deleted after 3C50''), \code{bytes}, \code{sha256} and \code{utc}; the NICER functions add
\code{tool}, \code{file_bytes}, \code{stop}, \code{rows} and offsets. At commit \texttt{666c2e0} it holds 50\,773 records
(\cref{tab:dlsummary}). Every input can be verified byte for byte after re-download. This was done once in full: on a
second machine (macOS), the 2\,689 v1--v2 raw files were re-downloaded from the logged URLs and all SHA256 values matched.
The rebuilt \file{features.npz} agreed with the original to $\le10^{-16}$.

\begin{table}[htbp]
\centering\small
\caption{Download log by host (\file{logs/downloads.jsonl}, commit \texttt{666c2e0}). Streamed NICER records count the
bytes read, although the bytes were not stored.}\label{tab:dlsummary}
\begin{tabularx}{\textwidth}{@{}l r r >{\raggedright\arraybackslash}X@{}}
\toprule
Host & Records & Volume & Tool(s) \\
\midrule
\file{heasarc.gsfc.nasa.gov} & 46\,731 & 5.2\,GB & \code{urllib} (\code{common.download}); \code{curl} for event files \\
\file{nasa-heasarc.s3.amazonaws.com} & 3\,384 & 229.5\,GB & \code{curl} range requests (prefix, tail, stream); \code{curl} in WSL (3C50) \\
\file{raw.githubusercontent.com} & 487 & 0.01\,GB & \code{urllib} \\
\file{maxi.riken.jp} & 156 & 0.06\,GB & \code{urllib} \\
\file{cdsarc.cds.unistra.fr} & 8 & $<0.01$\,GB & \code{urllib} \\
\file{arxiv.org} & 4 & $<0.01$\,GB & \code{urllib} \\
\file{ndownloader.figshare.com} & 3 & 0.06\,GB & \code{urllib} \\
\midrule
total & 50\,773 & 234.8\,GB & \\
\bottomrule
\end{tabularx}
\end{table}
